\chapter{Formal Languages and the Chomsky Hierarchy}

Part~I developed analytic combinatorics in a setting-agnostic way: a combinatorial class is simply a set of objects equipped with a notion of size, and the transfer theorems translate structural decompositions into functional equations for generating functions. This chapter grounds that framework in a concrete and historically fundamental family of classes, namely the \emph{formal languages}—subsets of the free monoid $\Sigma^*$ generated by grammars or recognized by machines. The structural organisation of such languages is the Chomsky hierarchy, a four-level stratification whose levels correspond, through the analytic-combinatorial lens, to progressively richer families of generating functions. The punchline, already foreshadowed in Chapter~5, is that regular languages have rational generating functions, unambiguous context-free languages have algebraic generating functions, and moving further up the hierarchy introduces genuinely transcendental and even exotic analytic behavior. Standard references for automata theory and formal languages are \cite{HopcroftUllman1979} and \cite{Sipser2013}.

\section{Alphabets, Strings, and Languages}

\begin{definition}
An \emph{alphabet} $\Sigma$ is a finite nonempty set whose elements are called \emph{symbols} or \emph{letters}. A \emph{string} over $\Sigma$ is a finite sequence $w = \sigma_1\sigma_2\cdots\sigma_k$ with each $\sigma_i \in \Sigma$. The \emph{length} of $w$ is $|w| = k$. The unique string of length zero is the \emph{empty string} $\varepsilon$.
\end{definition}

The set $\Sigma^*$ consists of all finite strings over $\Sigma$ including $\varepsilon$; it is the free monoid generated by $\Sigma$ under the operation of \emph{concatenation} $uv$ (juxtaposition). The subset $\Sigma^n \subset \Sigma^*$ collects strings of exactly length $n$, so $|\Sigma^n| = |\Sigma|^n$. We write $\Sigma^+ = \Sigma^* \setminus \{\varepsilon\}$ for the set of nonempty strings. A \emph{language} over $\Sigma$ is any subset $L \subseteq \Sigma^*$. The combinatorial quantity of primary interest is the sequence $(|L \cap \Sigma^n|)_{n \ge 0}$, which counts the words of each length; the generating function $L(z) = \sum_{n \ge 0} |L \cap \Sigma^n| z^n$ will be the central object of study.

\section{Deterministic Finite Automata and Regular Languages}

\begin{definition}
A \emph{deterministic finite automaton} (DFA) is a tuple $M = (Q, \Sigma, \delta, q_0, F)$ where $Q$ is a finite set of \emph{states}, $\delta : Q \times \Sigma \to Q$ is the \emph{transition function}, $q_0 \in Q$ is the \emph{start state}, and $F \subseteq Q$ is the set of \emph{accepting states}.
\end{definition}

The transition function extends to strings via $\hat{\delta}(q, \varepsilon) = q$ and $\hat{\delta}(q, w\sigma) = \delta(\hat{\delta}(q, w), \sigma)$ for $\sigma \in \Sigma$. The language \emph{accepted} by $M$ is $L(M) = \{w \in \Sigma^* : \hat{\delta}(q_0, w) \in F\}$.

\begin{example}
Let $\Sigma = \{a, b\}$ and $L = \{w \in \Sigma^* : w \text{ contains an even number of } a\text{'s}\}$. The DFA with $Q = \{q_0, q_1\}$, start state $q_0$, $F = \{q_0\}$, and transitions $\delta(q_i, a) = q_{1-i}$ and $\delta(q_i, b) = q_i$ accepts exactly $L$. The two states track the parity of the $a$-count.
\end{example}

A \emph{nondeterministic finite automaton} (NFA) relaxes $\delta$ to a relation $Q \times \Sigma \to 2^Q$ and allows $\varepsilon$-transitions. The subset construction converts any NFA with $n$ states into a DFA with at most $2^n$ states accepting the same language, so the two models are equivalent in expressive power.

\begin{definition}
A language $L$ is \emph{regular} if $L = L(M)$ for some DFA $M$.
\end{definition}

\section{The Myhill--Nerode Theorem}

The minimal DFA recognizing a language $L$ is characterised by an intrinsic equivalence relation on strings. Define $x \sim_L y$ if and only if for every $z \in \Sigma^*$, $xz \in L \iff yz \in L$: two strings are equivalent when they have identical sets of future continuations leading to acceptance. This is the \emph{Myhill--Nerode equivalence} of $L$.

\begin{theorem}[Myhill--Nerode]
A language $L$ is regular if and only if $\sim_L$ has finitely many equivalence classes. When $L$ is regular, the number of classes equals the number of states in the minimal DFA for $L$.
\end{theorem}

The proof in the forward direction is straightforward: two strings reaching the same state of a DFA are $\sim_L$-equivalent, so the number of classes is at most $|Q|$. Conversely, if there are only finitely many classes, one constructs a DFA whose states are the equivalence classes, and verifies it accepts $L$.

The theorem immediately shows that the language $\{a^n b^n : n \ge 0\}$ is not regular: the strings $\varepsilon, a, aa, aaa, \ldots$ are pairwise $\sim_L$-inequivalent, since $a^i z \in L \iff z = b^i$ while $a^j z \in L \iff z = b^j$, so distinct powers of $a$ have distinct future sets, giving infinitely many classes.

\section{Regular Expressions}

Regular expressions over $\Sigma$ are built by the following inductive syntax: the constants $\emptyset$ (empty language), $\varepsilon$ (language $\{\varepsilon\}$), and $\sigma$ (language $\{\sigma\}$) for each $\sigma \in \Sigma$ are regular expressions; if $R$ and $S$ are regular expressions then so are $R + S$ (union), $RS$ (concatenation), and $R^*$ (Kleene star, zero or more concatenations of strings from $R$). The language denoted by a regular expression is defined inductively from these operations.

\begin{theorem}[Kleene]
A language is regular if and only if it is denoted by some regular expression.
\end{theorem}

The equivalence is constructive in both directions: one converts a DFA to a regular expression by a state-elimination procedure, and one converts a regular expression to an NFA by structural induction. In the analytic-combinatorial vocabulary, the operations $+$, concatenation, and $*$ correspond exactly to disjoint union, the $\SEQ$ construction applied to a pair of classes, and $\SEQ$ applied to a single class without restriction on length; the close parallel between regular expressions and the symbolic method is visible here and will be developed further through weighted automata in Chapter~7.

\section{The Pumping Lemma for Regular Languages}

\begin{theorem}[Pumping Lemma, Regular]
If $L$ is a regular language, there exists a constant $p \ge 1$ (the \emph{pumping length}) such that for every $w \in L$ with $|w| \ge p$, there is a decomposition $w = xyz$ satisfying: $(i)$ $|xy| \le p$, $(ii)$ $|y| \ge 1$, and $(iii)$ $xy^k z \in L$ for every $k \ge 0$.
\end{theorem}

The proof is a direct application of the pigeonhole principle: let $p = |Q|$ where $Q$ is the state set of a DFA for $L$. Reading $w$ from $q_0$, the machine visits $|w|+1 \ge p+1$ states, so by pigeonhole some state $q$ is visited twice within the first $p+1$ steps; take $y$ to be the substring read between the two visits to $q$. Since the machine loops on $q$ when reading $y$, repeating or deleting $y$ does not change the final state, so $xy^kz \in L(M) = L$ for all $k \ge 0$.

To see that $L = \{a^n b^n : n \ge 0\}$ is not regular, suppose it were regular with pumping length $p$. Take $w = a^p b^p \in L$. The decomposition satisfies $|xy| \le p$, so $y = a^j$ for some $j \ge 1$; but then $xy^2 z = a^{p+j}b^p \notin L$, a contradiction.

\section{Context-Free Grammars and Pushdown Automata}

\begin{definition}
A \emph{context-free grammar} (CFG) is a tuple $G = (V, \Sigma, P, S)$ where $V$ is a finite set of \emph{nonterminal} symbols, $\Sigma$ is a finite alphabet of \emph{terminal} symbols (with $V \cap \Sigma = \emptyset$), $P \subseteq V \times (V \cup \Sigma)^*$ is a finite set of \emph{productions} written $A \to \alpha$, and $S \in V$ is the \emph{start symbol}.
\end{definition}

A \emph{derivation step} $\alpha A \beta \Rightarrow \alpha \gamma \beta$ applies a production $A \to \gamma$ to replace one occurrence of $A$. The reflexive-transitive closure $\Rightarrow^*$ defines reachability. The \emph{language generated} by $G$ is $L(G) = \{w \in \Sigma^* : S \Rightarrow^* w\}$. A language is \emph{context-free} (CFL) if it equals $L(G)$ for some CFG $G$.

\begin{example}
The grammar $S \to aSb \mid \varepsilon$ has a single nonterminal $S$, terminals $\{a, b\}$, and two productions. One verifies that $L(G) = \{a^n b^n : n \ge 0\}$: the only derivation of length $n$ produces $a^n b^n$, giving a generating function $L(z) = \sum_{n \ge 0} z^{2n} = 1/(1-z^2)$, which is rational here by a coincidence specific to this grammar.
\end{example}

\begin{example}
The grammar $S \to SS \mid (S) \mid \varepsilon$ generates the \emph{Dyck language} of balanced parentheses. The functional equation $S(z) = 1 + z^2 S(z)^2$ (using $z$ as a size marker for each matched pair) yields $S(z) = (1 - \sqrt{1-4z^2})/(2z^2)$, which is algebraic and not rational.
\end{example}

A grammar is \emph{ambiguous} if some string $w \in L(G)$ admits two distinct leftmost derivations (leftmost means always rewriting the leftmost nonterminal). Ambiguity is not merely a syntactic inconvenience: \emph{inherently ambiguous} languages exist for which no unambiguous grammar exists, and for such languages the symbolic method via the transfer matrix / Chomsky--Schützenberger approach fails to apply directly. The algebraic generating function correspondence from Chapter~5 requires an unambiguous grammar.

A \emph{pushdown automaton} (PDA) is a nondeterministic finite automaton augmented with a stack of unbounded depth; its transition function may push or pop symbols from the stack. The context-free languages are exactly the languages accepted by PDAs, in the same sense that DFAs capture the regular languages.

\begin{theorem}[Pumping Lemma, Context-Free]
If $L$ is a CFL with pumping length $p$, then every $w \in L$ with $|w| \ge p$ decomposes as $w = uvxyz$ with $|vxy| \le p$, $|vy| \ge 1$, and $uv^k x y^k z \in L$ for all $k \ge 0$.
\end{theorem}

The proof uses the parse tree of $w$: if $|w|$ exceeds $p = |\Sigma|^{|V|+1}$, the tree has height greater than $|V|$, so some nonterminal $A$ appears twice along a root-to-leaf path, giving a pumping structure in the tree. The language $\{a^n b^n c^n : n \ge 0\}$ is not context-free: taking $w = a^p b^p c^p$ and any decomposition $w = uvxyz$ with $|vxy| \le p$, the substring $vxy$ can span at most two of the three symbol types, so pumping distorts their equal counts.

\section{The Chomsky Hierarchy}

Formal grammars are classified by the form of their productions. This yields a strict four-level containment:

\begin{enumerate}
\item \textbf{Type 3 (Regular):} Productions of the form $A \to a$ or $A \to aB$ (right-linear grammars). Equivalent to finite automata. The class is closed under union, intersection, complement, concatenation, and Kleene star.

\item \textbf{Type 2 (Context-Free):} Productions of the form $A \to \alpha$ for $\alpha \in (V \cup \Sigma)^*$, with no restriction on $\alpha$. Equivalent to pushdown automata. The class is closed under union, concatenation, and Kleene star, but \emph{not} under intersection or complement: the intersection of two CFLs need not be context-free, as the example $\{a^n b^n c^k\} \cap \{a^k b^n c^n\} = \{a^n b^n c^n\}$ demonstrates.

\item \textbf{Type 1 (Context-Sensitive):} Productions of the form $\alpha A \beta \to \alpha \gamma \beta$ where $\gamma \neq \varepsilon$ (non-contracting). Equivalent to linear-bounded automata (Turing machines whose tape length is bounded by the input length). The language $\{a^n b^n c^n : n \ge 0\}$ is a canonical example, context-sensitive but not context-free.

\item \textbf{Type 0 (Recursively Enumerable):} Unrestricted grammars; productions may have any form. Equivalent to Turing machines. The \emph{halting language} $\{(\langle M \rangle, w) : M \text{ halts on } w\}$ is in this class but is not recursive (decidable), placing it strictly above the context-sensitive class.
\end{enumerate}

The inclusions $\text{Reg} \subsetneq \text{CFL} \subsetneq \text{CSL} \subsetneq \text{RE}$ are all strict, with canonical separating examples at each boundary. The regular and context-free levels are by far the most analytically tractable, and they account for nearly all structures arising in the study of language models.

\section{The Analytic-Combinatorial Payoff}

The structural information encoded in a formal language is reflected faithfully in the analytic type of its counting generating function $L(z) = \sum_{n \ge 0} |L \cap \Sigma^n| z^n$.

\begin{theorem}
If $L$ is a regular language, then $L(z)$ is a rational function of $z$.
\end{theorem}

The proof sketch is given in full in Chapter~7, but the idea is transparent: a DFA with state set $Q$ defines a linear system. For each state $q$, let $f_q(z) = \sum_{n \ge 0} |\{w \in \Sigma^n : \hat{\delta}(q_0, w) = q\}| z^n$ count the strings that land in $q$ after being read from the start state. The transitions $\delta(q, \sigma) = q'$ yield equations $f_{q'}(z) = \sum_{\sigma : \delta(q,\sigma) = q'} z \cdot f_q(z)$ (plus the contribution of $\varepsilon$ at the start state), a finite linear system over $\mathbb{Q}(z)$ whose solution is a vector of rational functions. The generating function $L(z) = \sum_{q \in F} f_q(z)$ is therefore rational.

\begin{theorem}[Chomsky--Schützenberger, Chapter 5]
If $L$ is generated by an unambiguous CFG, then $L(z)$ is algebraic over $\mathbb{Q}(z)$.
\end{theorem}

This is the content of Chapter~5's main theorem. The grammar's productions translate into a polynomial system for the generating functions of each nonterminal, and under the unambiguity hypothesis the solution is an algebraic function. Singularity analysis then yields the asymptotic form of $|L \cap \Sigma^n|$.

At level~1, the correspondence breaks down in a more severe way. Context-sensitive languages can exhibit generating functions with a \emph{natural boundary} on their circle of convergence—a circle beyond which no analytic continuation exists, so that the singularity analysis developed in Part~I simply does not apply. This is not a pathological edge case: for certain context-sensitive languages the coefficients grow faster than any geometrically bounded sequence, or oscillate in ways incompatible with the Darboux–transfer-matrix framework. The hierarchy thus presents a clean stratification:
\begin{align*}
\text{regular} &\;\longleftrightarrow\; \text{rational } L(z),\\
\text{unambiguous CFL} &\;\longleftrightarrow\; \text{algebraic } L(z),\\
\text{context-sensitive} &\;\longleftrightarrow\; \text{possibly transcendental or exotic } L(z).
\end{align*}
For the remainder of this book we work almost exclusively at the regular and context-free levels of the hierarchy (Types~3 and~2, respectively). The context-sensitive level appears in occasional warnings: if a model class tacitly encodes context-sensitive computations, one should not expect the clean rational or algebraic asymptotics that the transfer theorems deliver. Understanding precisely where a language sits in the hierarchy is therefore not merely a theoretical exercise—it determines which analytic tools are available.

Chapter~7 develops the theory of weighted finite automata, the probabilistic refinement of level~3 that assigns real-valued weights to transitions and thereby models probability distributions over $\Sigma^*$. This is the direct precursor to the $n$-gram language models of Part~III, and it is where the connection between regular languages and rational generating functions becomes computationally actionable.

\section*{Further Reading}

\begin{itemize}
  \item \textbf{J. E. Hopcroft and J. D. Ullman}, \emph{Introduction to Automata Theory, Languages, and Computation} (1979) \cite{HopcroftUllman1979} --- the definitive classical treatment of automata theory, formal languages, and the Chomsky hierarchy, with careful proofs of all major results including the Myhill--Nerode theorem and both pumping lemmas.
  \item \textbf{M. Sipser}, \emph{Introduction to the Theory of Computation} (2013) \cite{Sipser2013} --- a modern, accessible introduction to the theory of computation covering DFAs, NFAs, CFGs, pushdown automata, and undecidability; the standard graduate textbook for the material of this chapter.
  \item \textbf{P. Flajolet and R. Sedgewick}, \emph{Analytic Combinatorics} (2009) \cite{FlajoletSedgewick2009} --- Part~A develops the symbolic method and the connection between grammar structure and generating function type; the bridge between formal languages and analytic combinatorics that underlies this chapter's payoff section.
  \item \textbf{J. Berstel and C. Reutenauer}, \emph{Noncommutative Rational Series with Applications} (2011) \cite{BerstelReutenauer2011} --- treats rational and recognizable power series over noncommutative alphabets; the algebraic foundation for the transfer from regular languages to rational generating functions developed in Chapter~7.
\end{itemize}
